\section{An extension to nontrivial torus bundles}
\label{sec:twisted}
The Gowdy-to-Ernst examples on the three-torus with
matched full metrics use a periodic transformed $Q$,
so their scalar seed has zero-mean time velocity.
On a nontrivial torus bundle a nonzero current period is
admissible. This yields a rank-two collision for the
original four-field record with asymptotic speed
strictly between zero and one at every angle.

For an integer $n\ne0$ impose the type-II identification
\[
 (\theta,x,y)\sim(\theta+2\pi,x-ny,y),
 \qquad x,y\in\mathbb R/\mathbb Z.
\]
Take periodic $P,\lambda$ and
$Q(t,\theta+2\pi)=Q(t,\theta)+n$.
The orbit form $dx+Q\,dy$ descends under the displayed
identification. This is the standard unipotent
monodromy in twisted Gowdy \cite{rendall-twisted}.

\begin{proposition}
\label{prop:twisted}
For $n=1$ there are two distinct locally Gowdy vacuum
metrics with this compact torus-bundle topology, having
rank-two target maps on open sets and the same
$P(a),Q(a),Q_t(a),P(b)$ and initial metric anchor
in the common marked bundle coordinates.
Their asymptotic speeds are uniformly in $(0,1)$.
In the normalized current-one family generated from a
constant initial seed position, the terminal metric
potential distinguishes the two branches if the common
angular current period is nonzero.
\end{proposition}

\begin{proof}
Take the two geodesics from the subcritical row
of Table~\ref{tab:collisions}. Their speeds differ,
but they have common $q>0$ and common $P$ value
at both phase values $0$ and $\ell=3$.
Translate the initial $P$ to $P_0$ and rescale $Q$
by $e^{-P_0}$ to preserve the orthonormal components
of their initial velocities. Choose
$C=q e^{P_0}=1/(2\pi)$.
Use the phase
\[
 w(t,\theta)=\log(t/a)+\epsilon u_{k_0}(t)\cos(k_0\theta),
 \qquad u_{k_0}(a)=u_{k_0}(b)=0,\quad b/a=e^3 .
\]
Both seed metric potentials are periodic because
the initial orbit positions are constant, so the
periodic momentum constraint holds.
Apply \eqref{eq:ernst}. The transformed angular
period of $Q$ is
$C\int_0^{2\pi}t w_t\,d\theta=1$,
the required monodromy. The transformed metric potential
is periodic and the transformed $Q$ can be chosen common.
At both nodes $w_\theta=0$, and the initial
$Q_t$ and the measured $P$ values therefore agree.
The initial metric anchor can be taken common.
The missing initial velocities differ by a nonzero
multiple of $1+\epsilon\cos(k_0\theta)$.
The rank formula of Theorem~\ref{th:rank-two} still
holds at $a$, since $w_\theta(a)=0$.
Its factor $1+aP_w w_t$ is nonzero for both branches
for sufficiently small amplitudes: at zero amplitude
it is $1+g_j>0$.

The limit $t w_t\to1+\epsilon J_0(k_0a)\cos(k_0\theta)$
is positive everywhere for small $\epsilon$.
Then $w\to-\infty$ uniformly and $P_w\to-\sigma_j$
on each seed orbit. The transverse contribution
$t e^{\widetilde P}\widetilde Q_t$
vanishes as $t\downarrow0$.
The transformed speed is therefore
$1-\sigma_j(1+\epsilon J_0(k_0a)\cos(k_0\theta))$,
strictly within $(0,1)$ for both subcritical speeds
in the selected row and small enough amplitude.

For the final claim, normalize the seed current as
in Proposition~\ref{prop:ernst-fiber}.
The initial transformed $Q$ fixes the entire scalar
phase $\psi$, including the mean drift.
The energy comparison at the common angular origin
uses the integral $I$ in that proposition, positive
there whenever the common terminal phase is
nonzero. If terminal $P$ agrees, equality of terminal
$\lambda$ fixes the common $\kappa^2$ because the
$-2P-\log t$ terms in the transform agree.
Equality of terminal $P$ and \eqref{eq:ernst-fiber}
then determine $\zeta$ wherever $\psi(b)\ne0$.
For nonzero monodromy the mean terminal phase is
nonzero as $b>a$, so such an angle exists.
\end{proof}
